\section{MLReproMutate and Mutation Model}
\label{sec:mlrepromutate}

MLReproMutate is Python research software for applying controlled
mutations to reproducibility-relevant choices in machine-learning projects. It
evaluates whether an existing repository validation workflow detects the
resulting change.

The tool differs from model-level mutation approaches that perturb learned
parameters, neural-network structure, or input examples. Its mutation targets
are instead elements of the surrounding experimental software and
configuration: choices such as random seeds, dependency constraints, data
partitioning, and evaluation-protocol parameters.

The present study uses four implemented operators. Throughout the paper we use
the concise study labels \texttt{random-seed}, \texttt{dependency-pin},
\texttt{data-split}, and \texttt{cv-fold-count}. These labels identify the
mutation classes used in the empirical protocol; the corresponding
implementation classes have more explicit software-level names.


\subsection{Operator interface and mutation candidates}
\label{sec:mlrepromutate-interface}

MLReproMutate represents a mutation operator through two conceptually separate
operations: candidate detection and mutation application. An operator first
identifies syntactically supported locations in a repository and represents
each location as a mutation candidate. A candidate records the operator,
reproducibility category, target file, a human-readable description, and
operator-specific metadata needed to identify and validate the target.

Mutation application is then performed against a selected candidate rather than
by searching for a new target at execution time. The implementation verifies
that the candidate still corresponds to the expected source location before
rewriting the file. This design reduces the risk that a mutation is silently
applied to a different syntactic occurrence after candidate selection.

For Python-source operators, candidate detection uses the Python abstract syntax
tree (AST) and records source-position metadata. The mutation is applied only
after the relevant call and literal or expression have been matched again
against the stored candidate information. The dependency operator similarly
records the target requirements file, package, version, and line position and
verifies the expected pin before replacement.


\subsection{Reproducibility mutation operators}
\label{sec:mlrepromutate-operators}

Table~\ref{tab:mutation-operators} summarizes the four operators used in the
study. The transformations are intentionally narrow: each operator recognizes
a restricted syntactic form and applies one deterministic mutation.

\begin{table}[t]
\centering
\caption{Reproducibility-relevant mutation operators used in the empirical
study. Each transformation is deterministic once a candidate has been
selected.}
\label{tab:mutation-operators}

\vspace{0.45em}

\small
\renewcommand{\arraystretch}{1.30}
\setlength{\tabcolsep}{7pt}

\begin{tabularx}{\textwidth}{
    >{\raggedright\arraybackslash}p{0.18\textwidth}
    >{\raggedright\arraybackslash}X
    >{\raggedright\arraybackslash}p{0.31\textwidth}
}
\toprule
\addlinespace[2pt]

\textbf{Study label} &
\textbf{Supported target} &
\textbf{Transformation} \\

\addlinespace[2pt]
\midrule
\addlinespace[4pt]

\texttt{random-seed} &
Literal integer argument to
\texttt{random.seed},
\texttt{np.random.seed},
\texttt{numpy.random.seed}, or
\texttt{torch.manual\_seed} &
$N \rightarrow N+1$ \\

\addlinespace[5pt]

\texttt{dependency-pin} &
Exact dependency pin
\texttt{package==version}
in a supported requirements file &
\texttt{package==version}
$\rightarrow$
\texttt{package>=version} \\

\addlinespace[5pt]

\texttt{data-split} &
Explicit non-\texttt{None}
\texttt{stratify} argument in
\texttt{train\_test\_split} &
\texttt{stratify=<expr.>}
$\rightarrow$
\texttt{stratify=None} \\

\addlinespace[5pt]

\texttt{cv-fold-count} &
Explicit literal \texttt{n\_splits} in
\texttt{KFold}, \texttt{StratifiedKFold},
\texttt{RepeatedKFold}, or
\texttt{RepeatedStratifiedKFold} &
$N \rightarrow N+1$
for \texttt{n\_splits} \\

\addlinespace[4pt]
\bottomrule
\end{tabularx}

\vspace{0.3em}
\end{table}


\paragraph{\texttt{random-seed}.}
The random-seed operator targets supported seed-setting calls whose first
positional argument is an integer literal. It does not attempt to resolve
variables, expressions, configuration files, or dynamically supplied seed
values. For an eligible literal seed $N$, the mutation replaces the literal
with $N+1$. The transformation preserves the presence of explicit seeding while
changing the selected pseudorandom sequence.

\paragraph{\texttt{dependency-pin}.}
The dependency-pin operator recognizes exact requirement specifications of the
form \texttt{package==version} and relaxes the equality constraint to
\texttt{package>=version}. The textual mutation alone does not guarantee that a
different package version will be installed. For resolved-dependency
evaluation, MLReproMutate therefore constructs baseline and mutant dependency
environments separately and records the installed version of the target
distribution. If both specifications resolve to the same installed version,
the dependency mutation is treated as equivalent rather than as evidence that
the validation workflow failed to detect a dependency change.

\paragraph{\texttt{data-split}.}
The data-split operator targets supported scikit-learn
\texttt{train\_test\_split} calls containing an explicit
\texttt{stratify} keyword whose value is not \texttt{None}. The operator
replaces the complete stratification expression with \texttt{None}. Candidate
detection resolves supported import forms before accepting a call, reducing
false matches to unrelated functions with the same local name.

\paragraph{\texttt{cv-fold-count}.}
The cv-fold-count operator targets supported scikit-learn cross-validation
splitters with an explicit literal integer \texttt{n\_splits} argument. The
supported splitter classes are \texttt{KFold},
\texttt{StratifiedKFold}, \texttt{RepeatedKFold}, and
\texttt{RepeatedStratifiedKFold}. For an eligible fold count $N \geq 2$, the
operator replaces it with $N+1$.


\subsection{Baseline-first mutation evaluation}
\label{sec:mlrepromutate-evaluation}

MLReproMutate evaluates mutations relative to an executable unmodified
baseline. Before evaluating selected candidates, the validation workflow is run
against an isolated copy of the unmodified project. A baseline timeout or
non-zero exit status prevents the corresponding mutation from being interpreted
as a validation outcome.

After a successful baseline, each mutation is applied in an isolated project
workspace and the same validation command is executed against the mutated
project. At the engine level, a successful validation command corresponds to a
SURVIVED mutation, a non-zero exit status corresponds to a KILLED mutation, and
a validation timeout is represented separately.

The implementation also defines additional internal states such as INVALID and
EQUIVALENT. In particular, resolved dependency evaluation can distinguish a
syntactically changed dependency specification that resolves to the same
installed version from one that actually changes the resolved environment.

The empirical protocol in Section~\ref{sec:study-design} imposes stricter
study-level rules on top of these software-level execution states. Setup and
baseline failures are kept separate from mutation outcomes, and semantic
verification determines which evaluated mutations enter the confirmed
non-equivalent detection denominator.


\subsection{Isolation and auditability}
\label{sec:mlrepromutate-auditability}

Mutation evaluation is performed in temporary project sandboxes rather than by
modifying the source repository in place. Candidate metadata and mutation
results provide an explicit link between the selected source location, applied
transformation, validation execution, and resulting outcome.

This separation is important for empirical use. Candidate detection can be
performed independently from mutation execution, selected candidates can be
frozen before outcomes are observed, and the same mutation definition can be
reapplied to an isolated repository copy. MLReproMutate therefore provides the
mutation mechanism and execution primitives, while corpus construction,
workflow selection, semantic verification, and the study-specific stopping
rules remain explicit parts of the empirical protocol described in
Section~\ref{sec:study-design}.
